%! Parametric Functions and Rates of Change — Unit 8, Lesson 8.2 — Warm-Up (Answer Key)
\documentclass[10pt]{article}
\usepackage{precalculus-article}
\usepackage{precalculus-key}   % pulls in -boxes; do NOT also load -boxes
\newcommand{\TallMath}[1]{$\displaystyle #1\rule[-1.4em]{0pt}{3.2em}$}

\begin{document}

\pageheader{Unit 8, Lesson 8.2}{Warm-Up --- Answer Key}
\namedateperiod

\vspace{0.12in}

\begin{notesbox}{Spiral Review --- Key}

\textbf{1.}\quad Let $f(t) = 3t^2 - 12t$.

\vspace{4pt}
(a)\quad Is $f$ increasing or decreasing on the interval $[0, 2]$?
How do you know? (Hint: compare $f(0)$ and $f(2)$.)

\vspace{4pt}
$f(0) = 0$ and $f(2) = 3(4) - 24 = -12$. Since $f(2) < f(0)$, the function is
\ans{decreasing} on $[0, 2]$.

\vspace{0.10in}
(b)\quad Is $f$ increasing or decreasing on the interval $[2, 5]$?

\vspace{4pt}
$f(2) = -12$ and $f(5) = 75 - 60 = 15$. Since $f(5) > f(2)$, the function is
\ans{increasing} on $[2, 5]$.

\vspace{0.14in}
\textbf{2.}\quad Compute the average rate of change of $g(t) = t^2 - 4$
on the interval $[1, 4]$. Show your work.

\vspace{4pt}
$g(1) = 1 - 4 = -3$ and $g(4) = 16 - 4 = 12$.

$\text{ARC} = \dfrac{g(4) - g(1)}{4 - 1} = \dfrac{12 - (-3)}{3} = \dfrac{15}{3} =$ \ans{$5$}

\vspace{0.14in}
\textbf{3.}\quad Two particles both visit the point $(3, 5)$ in the plane.
Particle A arrives when $t = 1$; Particle B arrives when $t = 2$.
Explain why this is possible for parametric functions but is
\emph{not} possible for a function $y = f(x)$.

\vspace{4pt}
\ans{In a parametric function, two different values of the parameter $t$ can produce the same ordered pair $(x, y)$ at different times. The parameter $t$ tracks \emph{when} the particle is at a location, not just where. In a function $y = f(x)$, each input $x = 3$ produces exactly one output $y$, so only one particle can be at $x = 3$ with a given $y$-value --- the timing is not encoded.}

\begin{teachernote}
  Q1 previews the direction-of-motion analysis: students will use a similar
  comparison of values (rather than derivatives) to determine whether $x(t)$
  or $y(t)$ is increasing or decreasing. Q2 directly activates the average
  rate of change computation used in EK 4.3.A.4. Q3 surfaces the key
  conceptual advantage of parametric functions over $y = f(x)$: two different
  $t$-values can map to the same $(x, y)$ point, which EK 4.3.A.2 formalizes.
\end{teachernote}

\end{notesbox}

\end{document}
